\documentclass[11pt]{article}
\usepackage{mystyle}

\title{Rosen --- \S 8.2: Solving Linear Recurrence Relations\\ \large Selected Exercises}
\author{Alston}
\date{Semester 2, 2026}

\begin{document}
\maketitle

% ======================================================================
\begin{exercise}{1}
    Determine which of these are linear homogeneous recurrence
    relations with constant coefficients. Also, find the degree of
    those that are.
    \begin{enumerate}[label=(\alph*)]
        \item $a_n = 3a_{n-1} + 4a_{n-2} + 5a_{n-3}$
        \item $a_n = 2n\,a_{n-1} + a_{n-2}$
        \item $a_n = a_{n-1} + a_{n-4}$
        \item $a_n = a_{n-1} + 2$
        \item $a_n = a_{n-1}^2 + a_{n-2}$
        \item $a_n = a_{n-2}$
        \item $a_n = a_{n-1} + n$
    \end{enumerate}
\end{exercise}

% ======================================================================
\begin{exercise}{10}
    Prove Theorem 2.
\end{exercise}

% ======================================================================
\begin{exercise}{17}
    Prove this identity relating the Fibonacci numbers and the
    binomial coefficients:
    \[
    f_{n+1} = \binom{n}{0} + \binom{n-1}{1} + \cdots + \binom{n-k}{k},
    \]
    where $n$ is a positive integer and $k = \lfloor n/2 \rfloor$.
    [\emph{Hint}: Let
    $a_n = \binom{n}{0} + \binom{n-1}{1} + \cdots + \binom{n-k}{k}$.
    Show that the sequence $\{a_n\}$ satisfies the same recurrence
    relation and initial conditions satisfied by the sequence of
    Fibonacci numbers.]
\end{exercise}

% ======================================================================
\begin{exercise}{24}
    Consider the nonhomogeneous linear recurrence relation
    $a_n = 2a_{n-1} + 2^n$.
    \begin{enumerate}[label=(\alph*)]
        \item Show that $a_n = n2^n$ is a solution of this recurrence
        relation.
        \item Use Theorem 5 to find all solutions of this recurrence
        relation.
        \item Find the solution with $a_0 = 2$.
    \end{enumerate}
\end{exercise}

% ======================================================================
\begin{exercise}{36}
    Let $a_n$ be the sum of the first $n$ perfect squares, that is,
    $a_n = \sum_{k=1}^{n} k^2$. Show that the sequence $\{a_n\}$
    satisfies the linear nonhomogeneous recurrence relation
    $a_n = a_{n-1} + n^2$ and the initial condition $a_1 = 1$. Use
    Theorem 6 to determine a formula for $a_n$ by solving this
    recurrence relation.
\end{exercise}

% ======================================================================
\begin{exercise}{38}
    \begin{enumerate}[label=(\alph*)]
        \item Find the characteristic roots of the linear homogeneous
        recurrence relation $a_n = 2a_{n-1} - 2a_{n-2}$. [\emph{Note}:
        These are complex numbers.]
        \item Find the solution of the recurrence relation in part
        (a) with $a_0 = 1$ and $a_1 = 2$.
    \end{enumerate}
\end{exercise}

% ======================================================================
\begin{exercise}{40}
    Solve the simultaneous recurrence relations
    \[
    a_n = 3a_{n-1} + 2b_{n-1}, \qquad b_n = a_{n-1} + 2b_{n-1}
    \]
    with $a_0 = 1$ and $b_0 = 2$.
\end{exercise}

% ======================================================================
\begin{exercise}{41}
    \begin{enumerate}[label=(\alph*)]
        \item Use the formula found in Example 4 for $f_n$, the $n$th
        Fibonacci number, to show that $f_n$ is the integer closest
        to
        \[
        \frac{1}{\sqrt{5}} \left( \frac{1+\sqrt{5}}{2} \right)^n.
        \]
        \item Determine for which $n$, $f_n$ is greater than
        $\frac{1}{\sqrt{5}} \left( \frac{1+\sqrt{5}}{2} \right)^n$
        and for which $n$, $f_n$ is less than
        $\frac{1}{\sqrt{5}} \left( \frac{1+\sqrt{5}}{2} \right)^n$.
    \end{enumerate}
\end{exercise}

% ======================================================================
\begin{exercise}{44}
    (Linear algebra required) Let $A_n$ be the $n \times n$ matrix
    with 2s on its main diagonal, 1s in all positions next to a
    diagonal element, and 0s everywhere else. Find a recurrence
    relation for $d_n$, the determinant of $A_n$. Solve this
    recurrence relation to find a formula for $d_n$.
\end{exercise}

% ======================================================================
\begin{exercise}{45}
    Suppose that each pair of a genetically engineered species of
    rabbits left on an island produces two new pairs of rabbits at
    the age of 1 month and six new pairs of rabbits at the age of 2
    months and every month afterward. None of the rabbits ever die or
    leave the island.
    \begin{enumerate}[label=(\alph*)]
        \item Find a recurrence relation for the number of pairs of
        rabbits on the island $n$ months after one newborn pair is
        left on the island.
        \item By solving the recurrence relation in (a) determine the
        number of pairs of rabbits on the island $n$ months after one
        pair is left on the island.
    \end{enumerate}
\end{exercise}

% ======================================================================
\begin{exercise}{48}
    Show that the recurrence relation
    \[
    f(n) a_n = g(n) a_{n-1} + h(n),
    \]
    for $n \geq 1$, and with $a_0 = C$, can be reduced to a recurrence
    relation of the form
    \[
    b_n = b_{n-1} + Q(n) h(n),
    \]
    where $b_n = g(n+1) Q(n+1) a_n$, with
    \[
    Q(n) = \frac{f(1) f(2) \cdots f(n-1)}{g(1) g(2) \cdots g(n)}.
    \]
\end{exercise}

% ======================================================================
\begin{exercise}{50}
    It can be shown that $C_n$, the average number of comparisons
    made by the quick sort algorithm when sorting $n$ elements in
    random order, satisfies the recurrence relation
    \[
    C_n = n + 1 + \frac{2}{n} \sum_{k=0}^{n-1} C_k
    \]
    for $n = 1, 2, \dots$, with initial condition $C_0 = 0$.
    \begin{enumerate}[label=(\alph*)]
        \item Show that $\{C_n\}$ also satisfies the recurrence
        relation $nC_n = (n+1)C_{n-1} + 2n$ for $n = 1, 2, \dots$.
        \item Use Exercise 48 to solve the recurrence relation in
        part (a) to find an explicit formula for $C_n$.
    \end{enumerate}
\end{exercise}

% ======================================================================

\begin{exercise}{51}
    Prove Theorem 4.
\end{exercise}

% ======================================================================
\begin{exercise}{52}
    Prove Theorem 6.
\end{exercise}

% ======================================================================
\begin{exercise}{53}
    Solve the recurrence relation $T(n) = n\,T(n/2)^2$ with initial
    condition $T(1) = 6$ when $n = 2^k$ for some integer $k$.
    [\emph{Hint}: Let $n = 2^k$ and then make the substitution
    $a_k = \log T(2^k)$ to obtain a linear nonhomogeneous recurrence
    relation.]
\end{exercise}

\end{document}